В первой строке даны три целых числа $n$, $m$ и $k$~--- число тюремных камер, число вылазок Боузера и номер камеры, после которой расположен вход в подземелье ($1 \le n, m \le 1000$, $0 \le k \le n$). Если $k = 0$, то вход находится перед первой камерой.

Во второй строке даны $n$ целых чисел $a_i$~--- вместительности камер ($1 \le a_i \le 1000$).

Во третьей строке даны $m$ целых чисел $b_j$~--- число похищенных в $j$-ю вылазку принцесс ($1 \le b_j \le 1000$).